\documentclass[a4paper]{article}
\usepackage{amsfonts}
\usepackage{amsmath}
\usepackage{amssymb}
\usepackage{graphicx}

\begin{document}

\noindent \large Auteur: Sadia Nazary \\
\noindent \large Studierichting: Informatica \\
\noindent \large Oefeningenreeks 3A nr. 2.2 \\

\medskip

\normalsize

Bepaal de afgeleide van $f(x)=3x$ in de punten $a \in \{-2,0,3\}$ met behulp van de definitie \\

$f'(a)=\lim\limits_{h \to 0} {{f(a+h)-f(a)} \over h}$ \\

Voor $a=-2$ \\

$f(-2+h)=3(-2+h)=-6+3h$ \\

$f(-2)=-6$ \\

$f'(-2)=\lim\limits_{h \to 0} {{-6+3h-(-6)} \over h}$ \\

$=\lim\limits_{h \to 0} {{3h} \over h}$ \\

$=3$ \\

Voor $a=0$ \\

$f(0+h)=3h$ \\

$f(0)=0$ \\

$f'(0)=\lim\limits_{h \to 0} {{3h} \over h}$ \\

$=3$ \\

Voor $a=3$ \\

$f(3+h)=3(3+h)=9+3h$ \\

$f(3)=9$ \\

$f'(3)=\lim\limits_{h \to 0} {{9+3h-9} \over h}$ \\

$=\lim\limits_{h \to 0} {{3h} \over h}$ \\

$=3$ \\

$f'(a)=3$ voor alle $a \in \{-2,0,3\}$ \\

\begin{thebibliography}{999}
\end{thebibliography}

\end{document}